% @author Lluis Gesa (lluis.gesa@uab.cat)
%
% Enginyeria Del Software
% -----------------------------------------------------------------------------------



% definim quin estil general volem per el document
\documentclass[11pt]{article}

% indiquem quins paquets amb funcionalitat extra volem fer servir
\usepackage{a4wide, amsmath, tabularx, colortbl, fancyhdr, graphicx, lastpage, caption}
\usepackage[utf8]{inputenc}

% Indiquem que comencem a definir el document i els seus continguts.
\begin{document}

% Definim valors i propietats, creem macros que ens ajudaran a
% agilitzar coses repetides:

\newcommand{\NomProjecte}{Perdidos en la ignorancia}
\newcommand{\NomDocument}{Requerimentos}
\newcommand{\IdentificadorDocument}{GG-SL-N-REQ}
\newcommand{\VersioDocument}{v3.0}

% Definim una macro que cada cop que insertem \requeriment, creara una
% taula amb el contingut passat per paràmetre
\newcommand{\requeriment}[6]{
\begin{center}
	\begin{tabularx}{\linewidth}{|p{2.5cm}|X|}
	\hline
	\cellcolor[gray]{.8} \textbf{ID} & \cellcolor[gray]{.9} #1\\
	\hline
	\cellcolor[gray]{.8} \textbf{Title} & #2 \\
	\hline
	\cellcolor[gray]{.8} \textbf{Description} & #3\\
	\hline
	\cellcolor[gray]{.8} \textbf{Priority} & #4 \\
	\hline
	\cellcolor[gray]{.8} \textbf{Verification} & #5 \\
	\hline
	\cellcolor[gray]{.8} \textbf{Parents} & #6\\
	\hline
	\end{tabularx}
\end{center}
}

% Fent servir propietas del paquet fancyhdr per crear una capçalera
% que es mostrarà cada pàgina
\pagestyle{fancy}
\fancyhead[LO,LE]{\includegraphics[width = 3.5cm]{./imatges/uab-logo.png}}
\fancyhead[CO,CE]{\bf \NomProjecte \\ \vspace{\baselineskip}
  \NomDocument }
\fancyhead[RO,RE]{
{\bf Ref : \IdentificadorDocument} \\
Versió: \VersioDocument\\
Data: \today \\
Pag. \thepage ~of \pageref{LastPage}
}
\fancyfoot[LO,LE,CO,CE,RO,RE]{}
\textheight=19cm


%
% Creem la pàgina inicial :

% Insertem el nom del document en una taula centrada feng servir una
% font gran \huge

\begin{center}
\begin{tabular*}{15cm}{p{14.5cm}}
\\
\\
\\
\rowcolor[gray]{.9}\begin{center}\huge{\bf{\NomDocument}}\end{center}\\
\\
\end{tabular*}
\end{center}


%
% Insertem una taula que contindrà les abreviacions que hi hagi al
% nostre document
%
\small{%
\begin{center}
\begin{tabular}{|p{2cm}|p{4.75cm}|p{2cm}|p{4.75cm}|}
\hline
\multicolumn{4}{|>{\columncolor[gray]{.9}[2mm][2mm]}c|}{{\bf Abreviatures}} \\
\hline
UAB & Universitat Autonoma de Barcelona & ES & Enginyeria del Software \\
\hline
& & & \\
\hline
\end{tabular}
\end{center}
}

%
% Insertem una taula amb el historial de canvis del document
%
\small{%
\begin{center}
\begin{tabular}{|p{2cm}|p{2cm}|p{7.5cm}|p{2cm}|}
\hline
\multicolumn{4}{|>{\columncolor[gray]{.9}[2mm][2mm]}c|}{{\bf Historial de Revisions}} \\
\hline
{\bf Version} & {\bf Date} & {\bf Comments} & {\bf Autor} \\
\hline
1.0 & 08-Abril-2019 & Sprint 1 &
G43-2-02 \\
\hline
2.0 & 29-Abril-2019 & Sprint 2 & G43-2-02 \\
\hline
3.0 & 12-Maig-2019 & Sprint 3 & G43-2-02 \\
\hline
\end{tabular}
\end{center}
}




% Insertem en una nova pàgina, els indexs : de continguts, de figures
% i de taules, les comandes \tableofcontents ,  \listoffigures,
% \listoftables són gestionades per els paquets fets servir, i per
% defecte escriuen els titols de la secció en angles, per canviar-los,
% hem de 'renovar les comandes' :

\renewcommand{\contentsname}{Índex de continguts}
\renewcommand{\listfigurename}{Índex de figures}
\renewcommand{\listtablename}{Índex de taules}

% Ara li diem al LaTeX que inserti la informació
\newpage
\tableofcontents
\listoffigures
\listoftables

% Insertem les seccions que volem
\newpage

\section{Introducció}\label{sec:intro}

\begin{flushleft}
En este documento se pretende detallar de forma minuciosa el comportamiento deseado de todos los casos de uso del sistema. De esta forma, será mucho más fácil saber que debe hacer cada módulo cuando se quiera dividir el conjunto de tareas a programar.\\


Además, esto ayuda a anticiparse a posibles necesidades y problemas, lo que agiliza el proyecto y ayuda a determinar con mayor exactitud las herramientas a utilizar y tiempo a invertir.\\


En este mismo documento se proporcionará los diagramas de casos de uso, para detectar de forma sencilla los actores que intervienen y sus respectivas acciones.\\
\end{flushleft}

\section{Requeriments Funcionals}\label{sect:RF}


\requeriment{REQ-F-1-01}{USUARIO: Identificación de usuarios}{El usuario se conectará al sistema mediante un usuario y una contraseña.}{A}{R,T}{}


\requeriment{REQ-F-1-02}{ADMIN: creación y modificación de mapas}{
El administrador podrá crear nuevos mapas y modificar los existentes.}{A}{R,T}


\requeriment{REQ-F-1-03}{ADMIN: eliminar cuentas}{
El administrador podrá eliminar las cuentas de cualquier usuario.}{A}{R,T}


\requeriment{REQ-F-1-04}{ADMIN: crear cuentas}{
El administrador podrá crear nuevas cuentas para usuarios.}{A}{R,T}


\requeriment{REQ-F-1-05}{USUARIO: Empezar partida}{
El jugador podrá iniciar partidas cuando lo desee.}{A}{R,T}


\requeriment{REQ-F-1-06}{USUARIO: Borrar partida}{
El jugador podrá borrar sus partidas cuando lo desee.}{A}{R,T}


\requeriment{REQ-F-1-07}{ADMIN: Modificar contraseña}{
El administrador podrá modificar la contraseña del usuario.}{A}{R,T}


\requeriment{REQ-F-1-08}{ADMIN: Modificar ranking}{
El administrador podrá modificar el ranking de los usuarios.}{A}{R,T}


\requeriment{REQ-F-1-09}{USUARIO: Modificar o eliminar datos de la cuenta}{
El usuario puede modificar o eliminar sus datos de la cuenta.}{A}{R,T}


\requeriment{REQ-F-1-10}{SISTEMA: Gestionar juego}{
El sistema se encargará de gestionar todo lo relacionado con el juego (rankings, cuentas, etc.}{A}{R,T}


\requeriment{REQ-F-1-11}{SISTEMA: Gestionar niveles}{
El sistema gestionará los niveles de mayor a menos dificultad.}{A}{R,T}


\requeriment{REQ-F-1-12}{ADMIN: Ver estados de los usuarios}{
El administrador podrá ver el estado de los usuarios online (jugando, esperando, etc.).}{A}{R,T}


\requeriment{REQ-F-1-13}{USUARIO: Continuar partida}{
El usuario puede continuar la partida por donde la dejó.}{A}{R,T}


\requeriment{REQ-F-1-14}{USUARIO: Gestión de partidas existentes}{
El jugador tendrá en el menú la opción para gestionar sus partidas existentes.).}{A}{R,T}


\requeriment{REQ-F-1-15}{USUARIO: Nueva partida}{
El jugador tendrá en el menú la opción para empezar una partida nueva).}{A}{R,T}


\requeriment{REQ-F-1-16}{USUARIO: Visualización del ranking}{
El jugador tendrá en el menú la opción para ver los rankings.).}{A}{R,T}


\requeriment{REQ-F-1-17}{USUARIO: Modificación de los ajustes}{
El jugador tendrá en el menú la opción para modificar sus ajustes.).}{A}{R,T}


\requeriment{REQ-F-1-18}{USUARIO: Volver al menu principal}{
El jugador tendrá en el menú la opción para volver al menú principal.).}{A}{R,T}


\requeriment{REQ-F-1-19}{USUARIO: Salir del juego}{
El jugador tendrá en el menú la opción para salir del juego.).}{A}{R,T}


\requeriment{REQ-F-1-20}{USUARIO: Guardado manual}{
El jugador podrá guardar la partida manualmente.).}{A}{R,T}

\requeriment{REQ-F-1-20}{USUARIO: Guardado manual}{
El jugador podrá guardar la partida manualmente.).}{A}{R,T}

\requeriment{REQ-F-1-20}{USUARIO: Guardado manual}{
El jugador podrá guardar la partida manualmente.).}{A}{R,T}

\requeriment{REQ-F-1-21}{ADMIN: Borrar juego}{
El administrador podrá borrar el juego si quiere.).}{A}{R,T}

\section{Requeriments No Funcionals}\label{sect:RNF}


\requeriment{REQ-NF-1-01}{Escalabilidad del juego}{
El diseño debe ser escalable, es decir permitir ampliaciones.}{A}{R,T}{}

\requeriment{REQ-NF-1-02}{Diseño intutivo}{
El diseño del juego debe ser intuitivo y sencillo de manejar.}{A}{R,T}

\requeriment{REQ-NF-1-03}{Rendimiento}{
El sistema debe soportar como mínimo 300 jugadores.}{A}{R,T}

\requeriment{REQ-NF-1-04}{
Diseño adaptado a navegador web}{El diseño debe estar programado de forma que funcione desde un navegador.}{A}{R,T}

\requeriment{REQ-NF-1-05}{Seguridad}{
El sistema debe ser seguro (cifrados de contraseñas).}{A}{R,T}

\requeriment{REQ-NF-1-05}{Copias de seguridad}{
El servidor realizará copias de seguridad cada cierto tiempo para garantizar que no se pierden datos.}{A}{R,T}


\section{Traçabilitat}\label{sec:intro}

\begin{center}
\begin{tabular}{|c|c|c|}
\hline
{\cellcolor[gray]{.8} \bf Requeriment} & {\cellcolor[gray]{.8} \bf Cas d'ús} & {\cellcolor[gray]{.8} \bf Test}  \\
\hline
REQ-F-1-01 & CU-USUARI-01 & TP-U-01 \\
\hline
REQ-F-1-02 & CU-ADMIN-02 & TP-U-01 \\
\hline
REQ-F-1-03 & CU-ADMIN-03 & TP-U-01 \\
\hline
REQ-F-1-04 & CU-ADMIN-04 & TP-U-01 \\
\hline
REQ-F-1-05 & CU-USUARI-05 & TP-U-01 \\
\hline
REQ-F-1-06 & CU-USUARI-06 & TP-U-01 \\
\hline
REQ-F-1-07 & CU-ADMIN-07 & TP-U-01 \\
\hline
REQ-F-1-08 & CU-ADMIN-08 & TP-U-01 \\
\hline
REQ-F-1-09 & CU-USUARI-09 & TP-U-01 \\
\hline
REQ-F-1-10 & CU-SISTEMA-10 & TP-U-01 \\
\hline
REQ-F-1-11 & CU-SISTEMA-11 & TP-U-01 \\
\hline
REQ-F-1-12 & CU-ADMIN-12 & TP-U-01 \\
\hline
REQ-F-1-13 & CU-USUARI-13 & TP-U-01 \\
\hline
REQ-F-1-14 & CU-USUARI-14 & TP-U-01 \\
\hline
REQ-F-1-15 & CU-USUARI-15 & TP-U-01 \\
\hline

\end{tabular}
\captionof{table}{Matriu Traçabilitat}
\end{center}


% Finalitzem el document.
\end{document}
